% ============================================================
%  EXAMEN TEMA 7: TALLER INTEGRADOR Y EVALUACIÓN
%  Curso Introductorio — 3er Año
%  Duración: 90 minutos
%  Temas: Integración de lógica, álgebra, física y química
%         mediante problemas tipo Método de Pólya
%  Autor: Ing. Gabriel Astudillo
%  Nota: enunciado limpio (sin pistas ni desarrollo). Las
%  soluciones están en la "Clave de Respuestas" al final.
% ============================================================

\documentclass[12pt, letterpaper]{article}

% ── Paquetes ─────────────────────────────────────────────────

\usepackage{mathwizards-guia}
\mwguia{Examen — Tema 7: Taller Integrador y Evaluación}{Ing. Gabriel Astudillo $\cdot$ Curso 3er Año}

\begin{document}
% ============================================================

% ── Portada / Título ─────────────────────────────────────────
\mwportada{Examen del Tema 7}{Taller Integrador y Evaluación Final}{Problemas Complejos que Entrelazan Todas las Disciplinas}

\vspace{4pt}
\begin{cuadroinfo}
  \small
  \textbf{Nombre:} \rule{0.55\linewidth}{0.4pt} \quad \textbf{Fecha:} \rule{0.15\linewidth}{0.4pt} \\[6pt]
  \textbf{Tiempo disponible:} 90 minutos \quad$\bullet$\quad \textbf{Puntaje total:} 100 puntos \\[4pt]
  \textbf{Instrucciones:}
  \begin{enumerate}[leftmargin=*]
    \item Lee cada ejercicio con calma antes de resolverlo.
    \item Aplica explícitamente las cuatro fases del \textbf{Método de Pólya}: comprender, planificar, ejecutar y verificar.
    \item Muestra \emph{todos} tus pasos, incluidas las conversiones de unidades.
    \item \textbf{No olvides las unidades} y verifica la coherencia de tus resultados.
    \item No se permite el uso de calculadora.
    \item Escribe con letra clara y revisa tu examen antes de entregarlo.
  \end{enumerate}
\end{cuadroinfo}

\vspace{8pt}

% ============================================================
\section{Bloque I: Densidad, Conversiones y Geometría (30 puntos)}
% ============================================================

\begin{enumerate}[leftmargin=*]
  \item \textbf{(10 pts)} Un cilindro de aluminio ($d = 2{,}7\;\text{g/cm}^3$) tiene un radio de $4\;\text{cm}$ y una altura de $15\;\text{cm}$.
        \begin{enumerate}[label=\alph*)]
          \item Calcula su volumen en $\text{cm}^3$.
          \item Calcula su masa en gramos y en kilogramos.
          \item Si se sumerge completamente en agua, \textquestiondown cuántos mililitros de agua desplaza?
        \end{enumerate}

  \item \textbf{(10 pts)} Una esfera metálica tiene una masa de $1{,}95\;\text{kg}$ y una densidad de $7{,}87\;\text{g/cm}^3$.
        \begin{enumerate}[label=\alph*)]
          \item Convierte la masa a gramos.
          \item Calcula el volumen de la esfera en $\text{cm}^3$.
          \item Calcula su radio.
          \item \textquestiondown De qué metal podría tratarse?
        \end{enumerate}

  \item \textbf{(10 pts)} Un bloque tiene dimensiones $3\;\text{in} \times 2\;\text{in} \times 4\;\text{in}$ y una masa de $0{,}9\;\text{lb}$. Usa $1\;\text{in} = 2{,}54\;\text{cm}$ y $1\;\text{lb} = 453{,}6\;\text{g}$.
        \begin{enumerate}[label=\alph*)]
          \item Convierte las dimensiones a centímetros.
          \item Calcula el volumen en $\text{cm}^3$.
          \item Convierte la masa a gramos.
          \item Calcula la densidad en $\text{g/cm}^3$ e identifica el posible material.
        \end{enumerate}
\end{enumerate}

\newpage

% ============================================================
\section{Bloque II: Cinemática, Conversiones y Despeje (30 puntos)}
% ============================================================

\begin{enumerate}[leftmargin=*]
  \item \textbf{(10 pts)} Un auto viaja a $90\;\text{km/h}$.
        \begin{enumerate}[label=\alph*)]
          \item Convierte su velocidad a $\text{m/s}$.
          \item Si parte de $x_0 = 200\;\text{m}$, \textquestiondown cuál será su posición después de $3\;\text{min}$?
          \item \textquestiondown Cuánto tiempo (en segundos y en minutos) tardará en llegar a $x = 8\;\text{km}$?
        \end{enumerate}

  \item \textbf{(10 pts)} La velocidad del sonido en el agua es $1\,500\;\text{m/s}$ y en el aire es $340\;\text{m/s}$.
        \begin{enumerate}[label=\alph*)]
          \item Expresa ambas velocidades en $\text{km/h}$.
          \item \textquestiondown Cuántas veces más rápido viaja el sonido en el agua que en el aire?
          \item Si un buzo bajo el agua escucha una explosión ocurrida a $4{,}5\;\text{km}$ de distancia, \textquestiondown cuántos segundos tardó el sonido en llegarle?
        \end{enumerate}

  \item \textbf{(10 pts)} Despeja la variable indicada:
        \begin{enumerate}[label=\alph*)]
          \item $x = v_0 t + \frac{1}{2}at^2$ \quad (despeja $v_0$)
          \item $V = \frac{4}{3}\pi r^3$ \quad (despeja $r$)
          \item $p\,V = n\,R\,T$ \quad (despeja $T$)
        \end{enumerate}
\end{enumerate}

\newpage

% ============================================================
\section{Bloque III: Problemas Interdisciplinarios (25 puntos)}
% ============================================================

\begin{enumerate}[leftmargin=*]
  \item \textbf{(9 pts)} Un químico tiene $250\;\text{mL}$ de una sustancia líquida y mide su masa: $315\;\text{g}$.
        \begin{enumerate}[label=\alph*)]
          \item Calcula la densidad de la sustancia.
          \item \textquestiondown Esta sustancia flotaría en el agua ($d = 1\;\text{g/cm}^3$)? \textquestiondown Y en el mercurio ($d = 13{,}6\;\text{g/cm}^3$)?
          \item Usando $d_{\text{glicerina}} = 1{,}26\;\text{g/cm}^3$ y $d_{\text{aceite}} = 0{,}92\;\text{g/cm}^3$, \textquestiondown de qué sustancia podría tratarse?
        \end{enumerate}

  \item \textbf{(8 pts)} Una muestra metálica pesa $392\;\text{N}$ en la Tierra ($g = 9{,}8\;\text{m/s}^2$). Al sumergirla en una probeta con $2\;\text{L}$ de agua, el nivel sube a $7\;\text{L}$.
        \begin{enumerate}[label=\alph*)]
          \item Calcula la masa de la muestra.
          \item Calcula su volumen en litros y en $\text{cm}^3$.
          \item Calcula su densidad e indica si podría ser hierro ($7{,}87\;\text{g/cm}^3$) o cobre ($8{,}96\;\text{g/cm}^3$). Justifica.
        \end{enumerate}

  \item \textbf{(8 pts)} Se mezclan $500\;\text{g}$ de agua a $80\;^\circ\text{C}$ con $500\;\text{g}$ de agua a $20\;^\circ\text{C}$. Calcula la temperatura final de equilibrio. ($c_{\text{agua}} = 1\;\text{cal/(g}\cdot^\circ\text{C)}$.)
\end{enumerate}

\newpage

% ============================================================
\section{Conceptos Integradores (15 puntos)}
% ============================================================

\begin{enumerate}[leftmargin=*]
  \item \textbf{(5 pts)} Enumera y explica brevemente las cuatro fases del Método de Pólya.

  \item \textbf{(5 pts)} \textquestiondown Cuál es la diferencia entre \textbf{distancia recorrida} y \textbf{desplazamiento}? Da un ejemplo en el que la distancia sea mayor que el desplazamiento.

  \item \textbf{(5 pts)} \textquestiondown Cuál es la diferencia entre \textbf{calor} y \textbf{temperatura}?
\end{enumerate}

\newpage

% ============================================================
\ifmwclaves
  \section{Clave de Respuestas}
  % ============================================================

  \subsection*{Bloque I: Densidad, Conversiones y Geometría}

  \begin{enumerate}[leftmargin=*, label=\textbf{\arabic*.}]
    \item a) $V = \pi r^2 h = \pi(4)^2(15) = 240\pi \approx 754{,}0\;\text{cm}^3$.
          \quad b) $m = d\,V = 2{,}7 \times 754{,}0 \approx 2\,035{,}8\;\text{g} \approx 2{,}04\;\text{kg}$.
          \quad c) Desplaza $754{,}0\;\text{mL}$ de agua (igual a su volumen).

    \item a) $m = 1\,950\;\text{g}$.
          \quad b) $V = \dfrac{m}{d} = \dfrac{1\,950}{7{,}87} \approx 247{,}8\;\text{cm}^3$.
          \quad c) $V = \dfrac{4}{3}\pi r^3 \Rightarrow r = \sqrt[3]{\dfrac{3V}{4\pi}} = \sqrt[3]{\dfrac{3(247{,}8)}{4\pi}} \approx \sqrt[3]{59{,}15} \approx 3{,}9\;\text{cm}$.
          \quad d) La densidad $7{,}87\;\text{g/cm}^3$ corresponde al \textbf{hierro}.

    \item a) $3 \times 2{,}54 = 7{,}62\;\text{cm}$; \ $2 \times 2{,}54 = 5{,}08\;\text{cm}$; \ $4 \times 2{,}54 = 10{,}16\;\text{cm}$.
          \quad b) $V = 7{,}62 \times 5{,}08 \times 10{,}16 \approx 393{,}3\;\text{cm}^3$.
          \quad c) $m = 0{,}9 \times 453{,}6 \approx 408{,}2\;\text{g}$.
          \quad d) $d = \dfrac{408{,}2}{393{,}3} \approx 1{,}04\;\text{g/cm}^3$: podría ser un \textbf{plástico} (o material poco denso).
  \end{enumerate}

  \newpage

  \subsection*{Bloque II: Cinemática, Conversiones y Despeje}

  \begin{enumerate}[leftmargin=*, label=\textbf{\arabic*.}]
    \item a) $v = \dfrac{90}{3{,}6} = 25\;\text{m/s}$.
          \quad b) $3\;\text{min} = 180\;\text{s}$; $x_f = 200 + 25(180) = 200 + 4\,500 = 4\,700\;\text{m}$.
          \quad c) $8\;\text{km} = 8\,000\;\text{m}$; $t = \dfrac{8\,000 - 200}{25} = \dfrac{7\,800}{25} = 312\;\text{s} = 5\;\text{min}\;12\;\text{s}$.

    \item a) Agua: $1\,500 \times 3{,}6 = 5\,400\;\text{km/h}$. Aire: $340 \times 3{,}6 = 1\,224\;\text{km/h}$.
          \quad b) $\dfrac{1\,500}{340} \approx 4{,}41$ veces.
          \quad c) $4{,}5\;\text{km} = 4\,500\;\text{m}$; $t = \dfrac{4\,500}{1\,500} = 3\;\text{s}$.

    \item a) $v_0 = \dfrac{x - \frac{1}{2}at^2}{t} = \dfrac{2x - at^2}{2t}$.
          \quad b) $r = \sqrt[3]{\dfrac{3V}{4\pi}}$.
          \quad c) $T = \dfrac{p\,V}{n\,R}$.
  \end{enumerate}

  \newpage

  \subsection*{Bloque III: Problemas Interdisciplinarios}

  \begin{enumerate}[leftmargin=*, label=\textbf{\arabic*.}]
    \item a) $d = \dfrac{315}{250} = 1{,}26\;\text{g/cm}^3$.
          \quad b) Se \textbf{hunde} en el agua ($1{,}26 > 1$) y \textbf{flota} en el mercurio ($1{,}26 < 13{,}6$).
          \quad c) Coincide exactamente con la \textbf{glicerina} ($1{,}26\;\text{g/cm}^3$).

    \item a) $m = \dfrac{W}{g} = \dfrac{392}{9{,}8} = 40\;\text{kg}$.
          \quad b) $V = 7 - 2 = 5\;\text{L} = 5\,000\;\text{cm}^3$.
          \quad c) $d = \dfrac{40\,000\;\text{g}}{5\,000\;\text{cm}^3} = 8\;\text{g/cm}^3$. Está más cerca del \textbf{hierro} ($7{,}87$) que del cobre ($8{,}96$); podría tratarse de hierro o acero.

    \item $Q_{\text{cedido}} = Q_{\text{absorbido}}$: $500(80 - T_f) = 500(T_f - 20)$.
          $40\,000 - 500T_f = 500T_f - 10\,000 \Rightarrow 50\,000 = 1\,000T_f \Rightarrow T_f = 50\;^\circ\text{C}$.
          (Tiene sentido: como las masas son iguales, el equilibrio queda en el promedio de las temperaturas.)
  \end{enumerate}

  \newpage

  \subsection*{Conceptos Integradores}

  \begin{enumerate}[leftmargin=*, label=\textbf{\arabic*.}]
    \item \textbf{Comprender:} identificar datos, incógnitas y condiciones. \textbf{Planificar:} elegir una estrategia, fórmulas y conversiones. \textbf{Ejecutar:} desarrollar el plan paso a paso con unidades. \textbf{Verificar:} comprobar que el resultado y las unidades tienen sentido.

    \item La \textbf{distancia recorrida} es la longitud total del camino (escalar, siempre positiva). El \textbf{desplazamiento} es el cambio de posición (vectorial, puede ser negativo o cero). Ejemplo: caminar $3\;\text{m}$ al norte y luego $3\;\text{m}$ al sur: distancia $= 6\;\text{m}$, desplazamiento $= 0\;\text{m}$.

    \item La \textbf{temperatura} mide la agitación de las partículas de un cuerpo (en °C o K). El \textbf{calor} es la energía que se \emph{transfiere} entre dos cuerpos a diferente temperatura (en J o cal), siempre del caliente al frío.
  \end{enumerate}

\fi
\end{document}
